\subsection*{Conclusiones}

Tras analizar a fondo los datos y el desarrollo del proyecto, se determinaron las siguientes conclusiones:

\begin{itemize}
    \item \textbf{Reducción del tiempo de espera (Objetivo General):} La plataforma web con seguimiento GPS logró disminuir notablemente el tiempo de espera en los paraderos. Al comparar el pre-test y el post-test mediante la prueba estadística de Wilcoxon para las 20 observaciones pareadas, obtuvimos un valor $W = 0.0$ con un p-valor de $9.54 \times 10^{-7}$. Al ser este resultado menor al valor crítico de la tabla ($W_{\text{crit}} = 60$), rechazamos la hipótesis nula. El tiempo promedio que la gente espera pasó de $15.27$ minutos a apenas $8.05$ minutos, lo que significa un ahorro promedio de tiempo del $47.28\%$.
    
    \item \textbf{Situación de partida y diagnóstico (Objetivo Específico 1):} Los cuestionarios iniciales dejaron en claro que el desorden del transporte público en Arequipa estresa mucho a la población. El problema real no es solo que los buses tarden en pasar, sino que nadie sabe a qué hora llegarán. Esto justificó plenamente la creación del sistema de telemetría y estimación del arribo de las unidades.
    
    \item \textbf{Estructura técnica del software (Objetivo Específico 2):} La elección de Clean Architecture para programar la API del backend y el panel de administración fue acertada. Separar la lógica del negocio de los detalles de las bases de datos y la comunicación en tiempo real por WebSockets funcionó muy bien. Gracias a esto, el sistema es escalable y se le pueden agregar más funciones en el futuro sin comprometer lo que ya funciona.
    
    \item \textbf{Aceptación del usuario (Objetivo Específico 3):} Las encuestas finales después de usar el prototipo mostraron opiniones muy favorables. La interfaz en mapa con Leaflet agradó por su sencillez y facilidad de uso. A los participantes les pareció cómodo ver el bus avanzar en tiempo real, lo que demuestra la utilidad práctica de la solución.
    
    \item \textbf{Limitaciones detectadas:} Las pruebas se hicieron en un entorno de simulación basado en datos de la ruta Characato-Sabandia. Hay que tener en cuenta que en una implementación masiva real pueden surgir imprevistos como congestión por accidentes, desvíos de tráfico o zonas de Arequipa donde se pierda la señal móvil del GPS. Estas son variables del mundo real que podrían alterar la precisión de los tiempos estimados.
\end{itemize}

\newpage
\subsection*{Recomendaciones}

Para mejorar el sistema en futuras etapas del proyecto, sugerimos tomar en cuenta los siguientes puntos:

\begin{itemize}
    \item \textbf{Migración de motor de enrutamiento y Map Matching:} Aunque actualmente el sistema utiliza OSRM (Open Source Routing Machine) para proyectar las coordenadas sobre las calles y Leaflet junto a los mapas gratuitos de OpenStreetMap para la renderización visual en el cliente, se recomienda evaluar la migración a otros motores de código abierto como \textbf{Valhalla} o \textbf{GraphHopper}. Este cambio permitiría realizar búsquedas de rutas con múltiples criterios dinámicos y restricciones de giro de forma más flexible. Además, se aconseja añadir filtros de suavizado físico como el filtro de Kalman en el servidor para limpiar el ruido de las coordenadas GPS antes de dibujarlas en el mapa.
    
    \item \textbf{Integración de datos de tráfico y cierres viales en tiempo real:} Como el sistema se diseñó utilizando exclusivamente capas y servicios cartográficos gratuitos de OpenStreetMap, no cuenta con información de incidencias en tiempo real. Para resolver esto en un despliegue real, se recomienda incorporar APIs de datos de tráfico (como las de Google Maps, Mapbox o TomTom) o consumir feeds locales de incidencias viales. Así, el motor de enrutamiento podrá desviar las estimaciones y calcular el ETA considerando calles cerradas, accidentes o embotellamientos momentáneos.
    
    \item \textbf{Uso de algoritmos predictivos basados en machine learning:} Recomendamos descartar el cálculo lineal estático de velocidad-distancia para predecir el ETA (Tiempo Estimado de Arribo). En su lugar, sugerimos entrenar modelos de series temporales o redes neuronales recurrentes (LSTM) en el servidor de NestJS. Estos modelos deben entrenarse usando el historial de viajes acumulado en la base de datos de PostgreSQL, considerando factores clave como la hora del día, el día de la semana y el tráfico histórico de los cuellos de botella viales de la ciudad.
    
    \item \textbf{Hardware de rastreo GPS e infraestructura de comunicación:} Para la operación en producción con las empresas del SIT, conviene migrar de aplicaciones móviles en teléfonos inteligentes a dispositivos físicos dedicados (como rastreadores GPS de protocolos estándar tipo Coban o Queclink). Estos equipos deben enviar tramas directas a un servidor de sockets TCP/UDP o a un broker MQTT ligero, logrando estabilidad en la telemetría, menor consumo de datos y previniendo cortes en la transmisión por fallas de batería o cierres accidentales de app.
    
    \item \textbf{Notificaciones Push y Geofencing móvil:} Aconsejamos desarrollar notificaciones de alerta automáticas utilizando Progressive Web Apps (PWA) o Firebase Cloud Messaging. El sistema debe calcular en el backend cuándo un bus entra en un radio geográfico predefinido (geofence) cercano al paradero del usuario para enviar una vibración o alerta inmediata, liberando a la persona de la necesidad de mantener el navegador activo y la pantalla encendida para monitorear el mapa.
\end{itemize}
